\section{The bound for denominators}\label{sec:denominators}

This section applies Ribet's lemma to the Galois representations of the systems of Hecke eigenvalues in a space of forms and proves the bound for the denominator of the
projection to the lifts (\cref{thm:denominator-bound}). The auxiliary places of a triangle are added to the level; the tie parameters are realised by first-order
deformations of the norm maps from the deep auxiliary level to the base level; and the top coefficient of Ribet's element, an operator in the Hecke algebra, clears the
denominator.

\subsection{The Hecke algebra and tied traces}\label{ss:tied-traces}
Let base data $B$ fix a type of forms; the forms of level $L$ have principal level at auxiliary places $L_j$ and maximal level at the other places off the places of the data. A \emph{system} of level $L$ is an
isotypic space (forms with given Hecke eigenvalues at almost all places) containing a nonzero form of level $L$; it has a Galois representation $x.\mathrm{gal}$. The \emph{Hecke algebra} $\mathbb T$ is the ring of $\ell$-adic
functions $t$ on systems such that scaling each system $x$ by $\iota^{-1}(t(x))$ preserves the integral forms. Let $N_x$ be the twist of $x.\mathrm{gal}$ by a character $a$ trivial at the tame generators. The \emph{trace}
$\operatorname{tr}[P(h)g]$ is the function $x\mapsto\operatorname{tr}(P(N_xh)N_xg)$. For polynomials $H_j$ and variables $\alpha_j$, the \emph{tied ideal} $\mathfrak a\subseteq\mathbb T[\alpha]$ is generated by the $\ell^\nu\operatorname{tr}[g]$ and the traces
$\operatorname{tr}[(\gamma_j-1-\alpha_j(\varphi_j-1)H_j(\varphi_j))g]$ of the tied relations.
\lean{New.Patching.TaylorWilesPlaces.LocalChoice, New.Patching.TaylorWilesPlaces.nonempty_localChoice, New.Patching.TaylorWilesPlaces.LocalChoice.indLine, New.Patching.TaylorWilesPlaces.LocalChoice.indDet, New.Patching.TaylorWilesPlaces.LocalChoice.diamond, New.Patching.TaylorWilesPlaces.LocalChoice.transfer_indLine, New.Patching.TaylorWilesPlaces.LocalChoice.indDet_diamond, New.Patching.BaseLevel.placeRep, New.Patching.BaseLevel.lineRep, New.Patching.BaseLevel.commute_heckeAction_residueRep, New.Patching.BaseLevel.localFormRep_apply_eq_lineRep, New.Patching.BaseLevel.mem_levelSpace_proIwahori, New.Patching.BaseLevel.mem_parahoricEigenspace_line, New.Patching.BaseLevel.mem_parahoricEigenspace_det, New.Patching.BaseLevel.heckeEigenvalues_mem_heckeAlgebra, New.Patching.BaseLevel.IsGalRepTwist, New.Patching.BaseLevel.IsGalRepTwist.coe_apply, New.Patching.BaseLevel.charpoly_coeff_mem_heckeAlgebra, New.Patching.BaseLevel.IsCongruent, New.Patching.BaseLevel.exists_indicator_isCongruent, New.Patching.BaseLevel.traceElement, New.Patching.BaseLevel.tameTrace, New.Patching.BaseLevel.tameTrace_one, New.Patching.BaseLevel.tiedIdeal, New.Patching.BaseLevel.weightedSum_mul_heckeAction_eq_of_eigen, New.Patching.BaseLevel.tameTrace_smul_eq_diamond_sub_one, New.Patching.BaseLevel.aevalHecke, New.Patching.BaseLevel.tameTrace_smul_eq_aevalHecke, New.Patching.BaseLevel.weightedSum_indDet_mul_tameTrace, New.Patching.BaseLevel.norm_mul_traceElement_tameGenerator_eq_zero, New.Patching.BaseLevel.weightedSum_indLine_mem_levelSpace, New.Patching.BaseLevel.exists_aevalHecke_weightedSum, New.Patching.BaseLevel.map_aevalHecke_valMinAbs_mul_tameTrace, New.Patching.BaseLevel.lineRep_mul_eq_self, New.Patching.BaseLevel.heckeAction_mul_norm_mul_eq_zero}

At an auxiliary place $w=w_j$ of norm $q$, with $\ell^\nu\mid q-1$, a \emph{local choice} is a uniformiser $\varpi$ and a unit $z$ whose local Artin symbols are the classes of $\varphi_j$ and of $\gamma_j$, together with a logarithm
$\mathrm{ind}\colon\F_w^\times\to\Z/\ell^\nu$ with $\mathrm{ind}(\bar z)=1$; local choices exist by local class field theory, $\bar z$ generating $\F_w^\times$ modulo $\ell^\nu$-th powers. On $\GL_n(\F_w)$ we use $\mathrm{ind}\circ\det$, and on the stabiliser $P$ of the
line of the first basis vector the logarithm $\mathrm{ind}_1(g)=\mathrm{ind}(g_{11})$ of the first entry. A \emph{weighted norm} over a finite group $\Gamma$ acting on forms, for a homomorphism $\chi\colon\Gamma\to\Z/\ell^\nu$ and a weight $f\colon\Z/\ell^\nu\to\Z$,
is $\sum_{\gamma}f(\chi(\gamma))\gamma$. Two systems are \emph{congruent} if the characteristic polynomials of their representations $N$ agree modulo the maximal ideal at every Galois element.

\begin{proposition}[The dictionary at an auxiliary place]\label{prop:dictionary}
\begin{enumerate}[label=\textup{(\alph*)}]
\item The coefficients of the characteristic polynomials of the $N_xg$, and the indicator of the systems congruent to a given one, lie in $\mathbb T$.
\item At an auxiliary place $w$ off the base level, the tame generator acts as a diamond operator: on a form on which $\GL_n(\F_w)$ acts through $\psi\circ\mathrm{ind}\circ\det$, $\operatorname{tr}[(\gamma-1)gp(\varphi)]=(\langle\gamma\rangle-1)\operatorname{tr}[gp(\varphi)]$, with
$\langle\gamma\rangle=\mathrm{diag}(\bar z,1,\dots,1)$.
\item The Frobenius element acts through a normalised Hecke operator: if the stabiliser of a line acts through $\psi\circ\mathrm{ind}$ of the first entry with $\psi(n!)\ne1$, then $\operatorname{tr}[(\gamma-1)gp(\varphi)]=p(\tilde U)\operatorname{tr}[(\gamma-1)g]$,
$\tilde U=a(\varphi)U$ with $U$ the Hecke operator of $\mathrm{diag}(\varpi,1,\dots,1)$ for the pro-unipotent Iwahori subgroup.
\end{enumerate}
\end{proposition}

\begin{proof}
(a) \emph{Eigenvalues.} At a place $w$ of prime norm off the level $L$, the function sending a system to its $j$-th Hecke eigenvalue at $w$, read in $\bar\Q_\ell$, lies in $\mathbb T$: the Hecke operator is supported at a place not above $\ell$, an integral
combination of translations, and keeps the integral forms. \emph{Coefficients.} The twist by $a$ multiplies the $i$-th coefficient of the characteristic polynomial of $N_xg$ by a power of $a(g)$. The untwisted coefficients are limits, along
Frobenius elements near $g$ (\cref{thm:B-classfield}(d)), of coefficients of Hecke polynomials (\cref{ss:B-galois}), which are integers times eigenvalues; $\mathbb T$ is closed under such limits, the systems being finitely many and the
integral structure being preserved coefficientwise. \emph{The indicator.} For each system $x'$ not congruent to $\pi$ choose a coefficient $c_{x'}\in\mathbb T$ whose values at $x'$ and $\pi$ are incongruent, and put $s_\pi=\prod_{x'}(c_{x'}-c_{x'}(x'))$.
It vanishes at every non-congruent $x'$ and has norm at least $1$ at each $x$ congruent to $\pi$, since then $c_{x'}(x)\equiv c_{x'}(\pi)\not\equiv c_{x'}(x')$. The indicator is $1-\prod_{x}(1-s_\pi/s_\pi(x))$ over the congruent $x$, an integral polynomial in $s_\pi$ with
coefficients of norm at most $1$; this is the idempotent of one factor of the decomposition of a finite algebra over a henselian local ring.

(b) Let $y$ be a form of the system $x$ on which $\GL_n(\F_w)$ acts through $\psi\circ\mathrm{ind}\circ\det$. The maximal compact subgroup at $w$ acts on $y$ by $k\mapsto\psi(\mathrm{ind}(\det k))$, and the pro-$p$ Iwahori subgroup fixes $y$. By local--global
compatibility for Iwahori eigenvectors (\cref{thm:B-galois}(b)), inertia at $w$ acts on $N_x$ through the characters $\psi\circ\mathrm{ind}$ of the torus, all equal on the image of $\gamma$; so every eigenvalue of $N_x(\gamma)$ is $\iota(\psi(\mathrm{ind}\,\bar z))=\iota(\psi(1))$.
$N_x(\gamma)$ has finite order (inertia at $w\nmid\ell$ acts on $N_x$ through a finite quotient, $a$ being trivial at $\gamma$), hence is semisimple, hence the scalar $\iota\psi(1)$. On the other side the diamond $\langle\gamma\rangle$ acts on $y$ by
$\psi(\mathrm{ind}\det\langle\gamma\rangle)=\psi(1)$. Both sides of the claim, applied to $y$, are $(\psi(1)-1)\operatorname{tr}[gp(\varphi)]y$.

(c) On such a form the pro-unipotent Iwahori subgroup acts as the identity: an element of the stabiliser with first entry and determinant $1$ fixes the form, and the pro-unipotent Iwahori subgroup reduces into such elements.
The parahoric of the line acts by $k\mapsto\psi(\mathrm{ind}\,k_{11})$. By \cref{thm:B-galois}(b) the polynomial $\prod_i(X-\iota\xi_i(\varpi))$ of the Iwahori eigencharacter kills Frobenius on the Galois side, and the operator $U$ acts on the Iwahori-fixed vectors
of this eigencharacter with the eigenvalue attached to the first entry. Hence, in traces against an element $T$ of the Galois group, a polynomial in $\tilde U$ applied to $y$ acts as the same polynomial in the Frobenius matrix on the
summand where inertia acts by $\psi$; the factor $\gamma-1$ projects to this summand, as on the other summands inertia acts as the identity and $\psi(n!)\ne1$ separates them ($\psi$ is nontrivial on the image of $\mathrm{ind}$ of the elements of order dividing
$n!$). This gives the identity.
\end{proof}

\subsection{The family at the auxiliary places}\label{ss:aux-family}
Let $\mathrm{Nm}_j$ be the sum over $\GL_n(\F_{w_j})$, $n_{\rm Aut}$ the common multiple of the
stabiliser orders, and $\equiv$ equality after a reduction killing every operator $T$ with $\ell^\nu\mid2n!\,T$.
\lean{New.Selmer.exists_integral_combination_norm_ge, New.Selmer.exists_diagonal_mem_norm_ge, New.Selmer.pow_finrank_le_nnnorm_det, New.Selmer.norm_coeff_lt_one_of_mul_monic, New.Selmer.norm_eval_one_lt_one, New.Selmer.exists_integral_multiple_eval_one, New.Selmer.exists_one_le_nnnorm_one_sub_pow_mul_root, New.Selmer.exists_eventually_one_le_norm_eval_charpoly, New.Selmer.norm_coeff_comp_sub_charpoly_smul_lt_one, New.Selmer.normalizingAvatar_eq_of_isFrob, New.Selmer.isRoot_hecke_chain, New.Selmer.prod_X_sub_C_dvd_comp_of_hasHeckePolyAt, New.Selmer.eventually_exists_hasRootChain, New.Selmer.orthogonal_mem_invtSubmodule_heckePolyOnForms, New.Selmer.le_of_isotypic_mem}
\lean{New.Patching.BaseLevel.integralEnd_mul_apply, New.Patching.BaseLevel.isSMulRegular_integralEnd, New.Patching.BaseLevel.map_traceElement_sub_pow_mul, New.Patching.BaseLevel.exists_mul_traceElement_tameGenerator_eq_frob, New.Patching.BaseLevel.commute_of_mem_localOperators_of_ne, New.Patching.BaseLevel.exists_deformed_norm_family, New.Patching.BaseLevel.exists_norm_apply_eq_smul, New.Patching.BaseLevel.sub_dvd_pow_sub_smul_pow, New.Patching.BaseLevel.heckeAction_pow_apply, New.Patching.BaseLevel.norm_mul_heckeAction_apply, New.Patching.BaseLevel.sum_mul_residueRep_apply_eq_smul, New.Patching.BaseLevel.smul_top_coeff_mem_ellIntegralForms, New.Patching.BaseLevel.C_prod_sub_prod_mul_mem, New.Patching.BaseLevel.exists_heckeAlgebra_ribetElement, New.Patching.BaseLevel.exists_diagOp_mem_ellIntegralForms}

We first collect four facts. For a descent pair $O$ and an isotypic space $x$ write $u_x$ for the twist of $x.\mathrm{gal}$ by the normalising character $\chi$ of $O$ (\cref{ss:rep-type}), and $t=t_O$. A system is a \emph{weak lift} if every $u_x(g)$ has the eigenvalue $1$.
A system has a \emph{root chain} at $w$ for $H$ if its Hecke polynomial, with the variable rescaled by the Frobenius value of $\chi$, is $\prod_{i<t}(X-q^i)$ times a divisor of $H$.

\begin{lemma}[Chains of roots]\label{lem:root-chains}
\begin{enumerate}[label=\textup{(\alph*)}]
\item Let $l_{pk}\in\bar\Q_\ell$. There are integral $d_p$ with $\norm{\sum_pl_{pk}d_p}\ge\norm{l_{kk}}$ for every $k$. Consequently, if each basis vector $b_k$ has an operator in a set closed under integral combinations, diagonal in the basis $b$ with
eigenvalues $l_k$, then one operator in the set, diagonal in $b$, has eigenvalues $\mu_k$ with $\norm{\mu_k}\ge\norm{l_{kk}}$; and if $\norm{c\,\iota\mu}\ge\varepsilon$ for every eigenvalue $\mu$ of an operator on a space of dimension $N$, then
$\norm{c^N\iota\det}\ge\varepsilon^N$.
\item If $f\ne0$ lies in a weak lift $x$ and has Hecke polynomial $S$ at a split $w\nmid\ell$ of norm $q$, then $X_0q^i$, $i<t$, are roots of $S$, $X_0$ the Frobenius value of the character of $W$ (the inverse of $\chi$ at a geometric Frobenius);
so $\prod_{i<t}(X-q^i)$ divides $S(X_0X)$.
\item Let $u_x$ have integral characteristic polynomials with the root $1$. For most Frobenius elements $\varphi$ at places $w$: if $\det(X-u_x(\varphi))\equiv S'\prod_{i<t}(X-\mathrm{cyc}(\varphi)^{-i})$ modulo the maximal ideal with $\norm{S'(1)}\ge1$, then $x$ has a root chain at $w$
for an integral $e$ with $e(1)=1$.
\item The orthogonal complement of a space $Z$ stable under the Hecke polynomials off the level is stable under them; and $Z\subseteq N'$ if every nonzero form of $Z$ lying in an isotypic space lies in $N'$.
\end{enumerate}
\end{lemma}

\begin{proof}
(a) The linear forms $F_k(X)=\sum_pl_{pk}X_p$ over the infinite residue field of $\bar\Q_\ell$: after scaling, finitely many nonzero hyperplanes do not cover the residue space, so some integral point $d$ has $\norm{F_k(d)}$ equal to the largest norm of a
coefficient of $F_k$, at least $\norm{l_{kk}}$, for every $k$ (forms that vanish identically are handled by the trivial case $l_{kk}=0$). The combination $\sum_pd_p\Lambda_p$ of the given diagonal operators is diagonal with eigenvalues $\sum_pl_{pk}d_p$. The
determinant is the product of the eigenvalues with multiplicity, $N$ in number.

(b) A geometric Frobenius $s^{-1}$ at $w$ has the eigenvalue $1$ on $u_x$, so $\chi(s^{-1})^{-1}$ is an eigenvalue of $x.\mathrm{gal}(s^{-1})$, which is killed by the Hecke polynomial (\cref{thm:B-galois}(a)); this makes $X_0$ a root of $S$. The form
lies in a unitary representation and $\abs{X_0}^2=q^{r}$, so by the chain lemma (\cref{thm:B-chains}(a)), with $N=n$ and $k=t$, $X_0,X_0q,\dots,X_0q^{t-1}$ are roots.

(c) For most Frobenius elements the place is split, prime to $\ell$ and off the level, and the characteristic polynomial of $x.\mathrm{gal}(\varphi)$ is $\iota$ of the Hecke polynomial. Let $B=\prod_{i<t}(X-q^i)$, which divides the rescaled Hecke polynomial
by (b), with quotient $T$. As $\mathrm{cyc}(\varphi)^{-1}=q$, $\iota(T)\iota(B)$ is congruent to $\chi^n$ times the characteristic polynomial of $u_x(\varphi)$, hence to $S'B$. Monic factors cancel in bounds on coefficients: if $B$ is monic and $(T-S')B$ has all
coefficients of norm below one, so has $T-S'$. Hence $T$ is integral and $T(1)\equiv S'(1)$ has norm at least $1$, and $e=T/T(1)$ is integral with $e(1)=1$ and $T\mid e$.

(d) The adjoint of the Hecke polynomial at $w$ is a Hecke polynomial at $\bar w$ (\cref{prop:point-forms}(b)), which preserves $Z$. For the second claim view the forms as smooth forms: $Z$ is spherical off the chosen places and stable under the Hecke
operators off them and their conjugates, so it is spanned by joint eigenforms (\cref{ss:B-hecke-transfer}); these lie in isotypic spaces, hence in $N'$.
\end{proof}

\begin{theorem}[The deformed norm family]\label{thm:deformed-norm}
\begin{enumerate}[label=\textup{(\alph*)}]
\item (Chains for weak lifts.) If every element has eigenvalue $1$ on the normalised representation of a system $x$, its Hecke polynomial at a split $w\nmid\ell$ has the roots $X_0q^i$, $i<t$; and for most Frobenius
elements, if $\det(X-u(\varphi))\equiv S\prod_{i<t}(X-\mathrm{cyc}(\varphi)^{-i})$ modulo the maximal ideal with $\abs{S(1)}\ge1$, the rescaled Hecke polynomial is a root chain times a divisor of an integral polynomial with value $1$ at $1$.
\item (First-order term.) If the base systems where $\delta$ does not vanish have a root chain at $w_j$, there is an integral operator $\partial$ made of translates at $w_j$ with $\mathrm{Nm}_j\operatorname{tr}[(\gamma-1)g]=0$ and
$\partial\operatorname{tr}[(\gamma-1)g]\equiv\mathrm{Nm}_j\operatorname{tr}[(\varphi-1)H(\varphi)g]$, and $n_{\rm Aut}\partial(g-1)\equiv\mathrm{Nm}_jZ$ with $\delta\,\mathrm{Nm}\,Z=0$.
\item (The family.) With $\delta=n_{\rm Aut}^2\delta'$ and root chains at all $w_j$, the coefficients $\mathcal N_S$ of $\prod_j(\mathrm{Nm}_j+\alpha_j\partial_j)$ are integral operators that kill $\mathfrak a$ and, up to powers of $\delta$, factor through the
base level.
\item (The top coefficient clears denominators.) Let $\mathcal D\in\mathbb T[\alpha]$ satisfy $(\delta_0-n_0)\mathcal D\in\mathfrak a$, let $l$ be a form of base level with $dl$ integral, and $y$ with $A\delta_0l=n_0Ay$, $\delta_0y=n_0y$ and $\mathcal D_Sy=0$ for
$S\ne J$. If $2n!(n_0n_{\rm Aut}^2)^{\#J+1}$ divides $\ell^\nu$, then $dn_{\rm Aut}A^{\#J+1}$ times the top coefficient $\mathcal D_J$ applied to $y$ is integral.
\end{enumerate}
\end{theorem}

\begin{proof}
(a) is \cref{lem:root-chains}(b), (c).

(b) \emph{The operator.} Fix a local choice at $w=w_j$ and write $t=t'+1$, $p=(X-q^{t'})H$. Let $\rho$ be the action of $\GL_n(\F_w)$, $P$ the stabiliser of the line, $L$ the weighted norm over $P$ with the weight $\nu\circ\mathrm{ind}_1$ ($\nu$ the lift of
least absolute value), and $B(f)=p(\tilde U)L_f$ for the weighted norms $L_f$ over $P$ with weight $f\circ\mathrm{ind}_1$; $B$ is additive in $f$, with values integral operators made of translates at $w$, as $U$ and the translations are. Put
$\partial=-(\sum_hh)B(\nu)(\sum_hh^{-1})$, the sums over fixed representatives $h$ of $\GL_n(\F_w)/P$ (\cref{ss:norms-logs}). \emph{The norm kills the tame trace.} By \cref{prop:dictionary}(b) with $p=1$ and the weight $1$, $\mathrm{Nm}_j\operatorname{tr}[(\gamma-1)g]=(\langle\gamma\rangle-1)\mathrm{Nm}_j\operatorname{tr}[g]$,
and $\langle\gamma\rangle$ fixes the norm. \emph{The tie.} Let $\sigma$ be the reduction killing every $T$ with $\ell^\nu\mid2n!T$. On eigenforms of $P$ for the characters $\psi\circ\mathrm{ind}_1$ with $\psi(n!)\ne1$, \cref{prop:dictionary}(c) gives
$p(\tilde U)\operatorname{tr}[(\gamma-1)g]=\operatorname{tr}[(\gamma-1)gp(\varphi)]$. The difference $\mathcal E(f)$ of the two sides after the weighted norm $L_f$ is additive in $f$ and vanishes at the weights $f(\cdot-n!)-f$: the values of these weighted norms lie in the
span of the eigenforms with $\psi(n!)\ne1$, as the characters form a basis of the functions on $\Z/\ell^\nu$ and the shift by $n!$ multiplies $\psi$ by $\psi(-n!)$. By \cref{lem:norms-logs}(e) with $k=n!$, $\sigma(\mathcal E(\nu))=0$. Then
$\sigma(\partial\operatorname{tr}[(\gamma-1)g])=\sigma(-(\sum_hh)L\operatorname{tr}[(\gamma-1)gp(\varphi)](\sum_hh^{-1}))$. By the transfer of the logarithm of the first entry (\cref{lem:norms-logs}(c), (d)), $(\sum_hh)L(\sum_hh^{-1})\equiv$ the logarithm of $\mathrm{ind}\circ\det$; by
\cref{prop:dictionary}(b) on the weighted norm with weight $\nu\circ\mathrm{ind}\circ\det$ the factor $\gamma-1$ becomes $\langle\gamma\rangle-1$, and $(\langle\gamma\rangle-1)L_{\mathrm{ind}\det}\equiv-\mathrm{Nm}_j$ by \cref{lem:norms-logs}(d), as $\mathrm{ind}\det\langle\gamma\rangle=1$. Hence
$\partial\operatorname{tr}[(\gamma-1)g]\equiv\mathrm{Nm}_j\operatorname{tr}[gp(\varphi)]$, and $p(\varphi)=(\varphi-q^{t'})H(\varphi)$ may be replaced by $(\varphi-1)H(\varphi)$, as $q\equiv1$ modulo $\ell^\nu$ and the two polynomials differ by a multiple of $q^{t'}-1$. \emph{The norm condition.}
$P$ fixes the values of $B(1)$: $U$ preserves the forms fixed by the parahoric subgroup, and on forms an element of $P$ acts through any lift to the parahoric. By \cref{lem:norms-logs}(f), with $L$ in the role of the invariant functional,
$n_{\rm Aut}\partial(g-1)\equiv\mathrm{Nm}_jZ$ with $Z=-\Psi B(1)Q$. \emph{$\delta\mathrm{Nm}Z=0$.} The values of $\mathrm{Nm}\,B(1)$ have base level. Let $x$ be a base system with $\delta(x)\ne0$ and $w'$ a form of base level in $x$; it has a root chain at $w_j$, so by the
chain lemma (\cref{thm:B-chains}(b)) it is orthogonal to the values of $B(1)=((X-q^{t'})H)(\tilde U)$ applied to forms fixed by the pro-$p$ Iwahori subgroup, and the finite groups fix $w'$. Hence $\delta\,\mathrm{Nm}\,B(1)y$ is orthogonal to all
$\delta w'$ with $w'$ of base level, among them itself; it is zero. Hence $\delta\,\mathrm{Nm}\,Z=0$.

(c) \emph{The coefficients.} Operators made of translates at different auxiliary places commute, so the $\mathcal A_j$ generated by the actions at $w_j$ are pairwise commuting, and \cref{lem:norms-logs}(a) gives the coefficients $\mathcal N_S$ of
$\prod_j(\mathrm{Nm}_j+\alpha_j\partial_j)$, with $\mathcal N_\emptyset=\mathrm{Nm}$; they commute with $\mathbb T$, which commutes with the residue actions at places off the base level. \emph{Killing $\mathfrak a$.} It suffices to treat multiples $G'z$ of the generators $z$ of $\mathfrak a$.
For $z=\ell^\nu\operatorname{tr}[g]$ the claim holds after $\sigma$. For a tied relation $z=\operatorname{tr}[(\gamma_j-1-\alpha_j(\varphi_j-1)H_j(\varphi_j))g]$, write $t=\operatorname{tr}[(\gamma_j-1)g]$ and $r=\operatorname{tr}[(\varphi_j-1)H_j(\varphi_j)g]$, so $z=t-\alpha_jr$; by (b), $\mathrm{Nm}_jt=0$ and
$\partial_jt\equiv\mathrm{Nm}_jr$, and \cref{lem:norms-logs}(a) gives that the top coefficient of $\mathcal N\cdot G'(t-\alpha_jr)$ is $\equiv0$. As $\mathfrak a$ is an ideal, this covers all square-free monomials. \emph{The factorisation.} By Higman's criterion
there is an integral $\Psi$ with $\mathrm{Nm}(\Psi y)=n_{\rm Aut}y$ for forms of base level; with (b) the hypotheses of \cref{lem:norms-logs}(b) hold, so $\delta^{\#J}\mathcal N_S\equiv Y_S\mathrm{Nm}$ with integral $Y_S$ and $Y_\emptyset=\delta^{\#J}$.

(d) Put $\delta=n_{\rm Aut}^2A\delta_0$, $n_1=n_0n_{\rm Aut}^2$ and $i_0=\#J+1=i+1$. For the base systems where $A\delta_0$ does not vanish the root chains hold, so (c) applies with this $\delta$. Let $G$ be the coefficient of $\prod_j\alpha_j$ in $\mathcal N\mathcal D$. As
$(\delta_0-n_0)\mathcal D\in\mathfrak a$ and $\delta_0-n_0$ divides $(n_{\rm Aut}^2A\delta_0)^k-(n_1)^kA^k$, modulo $\ell^\nu/(2n!)$, a multiple of $n_1^{i_0}$,
\[
  n_1^{i_0}A^{i_0}G\equiv\delta^{i_0}G=(\delta^i\mathcal N)(\delta\mathcal D)\equiv\sum_SY_{J\setminus S}\mathrm{Nm}\,\delta\,\mathcal D_S .
\]
Apply both sides to the integral form $u=\Psi(dl)$, whose norm is $dn_{\rm Aut}l$. With $A\delta_0l=n_0Ay$, $\mathrm{Nm}\,\delta\,s\,u=n_1dn_{\rm Aut}As\,y$ for $s\in\mathbb T$, and the terms with $S\ne J$ vanish as $\mathcal D_Sy=0$; with $Y_\emptyset=\delta^i$, which acts on $y$ as
$(n_1A)^i$ since $\delta_0y=n_0y$, only $n_1^{i_0}dn_{\rm Aut}A^{i_0}\mathcal D_Jy$ remains. The left side applied to $u$ is integral, and cancelling $n_1^{i_0}$, allowed as $2n!\,n_1^{i_0}\mid\ell^\nu$, gives the claim.
\end{proof}

\subsection{Ribet's element and the clearing bound}\label{ss:ribet-element}
For $\ell$ and the data put $N=n_{\rm ch}$, the order of $\prod_w\GL_{\mathrm r\mathrm t}(\F_w)$ over the chosen places. The \emph{odd class bound} is a positive integer at most $2^{4(N_E+\mathrm t+3)^{27}}h_E$ whose valuation bounds the lengths of the odd
classes in $(O/\ell^\nu)(-e)$, $0<e<t$, and of the unramified homomorphisms for $e=0$. The \emph{local defect bound} is $\mathrm{loc}_\ell=4f\,\mathrm{defect}_\ell$ (\cref{sec:local-loss}), the \emph{twist error} is $e'_\ell=(4f+3)\mathrm r\,v_\ell(N)+2^{-\ell}$,
the \emph{denominator error} is
\[
  e_\ell=\mathrm t\,v_\ell(X)+8f\mathrm t\,v_\ell(N)+2v_\ell(N)+\bigl((4\mathrm t+1)h_\Delta^{22}+h_\Delta\bigr)v_\ell(N),\qquad h_\Delta=\mathrm r^8,
\]
$X$ the odd class bound, and the \emph{clearing bound} is
\[
  \beta_\ell=\mathrm{loc}_\ell+(1+2\mathrm t)e'_\ell+e_\ell+v_\ell(N).
\]
A twist $\psi$ of modulus $Q$ is \emph{auxiliary} with ramification at most $\mathrm{ram}$ if $Q$ is auxiliary for the bad places, $\psi$ comes from the norm-one torus, and the valuations of $\psi(\sigma_P)-1$ at the tame generators
sum to at most $\mathrm{ram}$.
\lean{New.Selmer.oddClassBound, New.Selmer.oddClassBound_ne_zero, New.Selmer.oddClassBound_le, New.Selmer.eventually_length_shaOdd_le, New.Selmer.localDefectBound, New.Selmer.denominatorError, New.Selmer.denominatorError_eq_natCast, New.Selmer.denominatorError_ne_top, New.Selmer.clearingBound, New.Selmer.add_denominatorError_le_mul_clearingBound, New.Selmer.mul_padicValNat_add_le_add_denominatorError, New.Selmer.ribetCocycles.apply_eq_zero, New.Selmer.ribetCocycles.mem_selmerModule, New.Selmer.isTraceIdeal_withConjInv, New.Selmer.restrictToCoeffRing, New.Selmer.IsTraceIdeal.of_isLatticeIn, New.Selmer.summandRep_some_apply, New.Selmer.summandCoordIndexEquiv, New.Selmer.blockDiagonal'_summandRep, New.Selmer.exists_blockDiagonal_lattice, New.Selmer.exists_blockDiagonal_lattices, New.Selmer.IsTwistedLattice.twistRep_apply_eq_of_eq_one, New.Selmer.IsTwistedLattice.twistRep_eq_one_of_mem_inertiaAbove, New.Selmer.length_ribetCocyclesVanishingOn_le, New.Selmer.hasTriangleCost_of_isCharacterCocycle, New.Selmer.isTameType_of_mem_placesAboveM, New.Selmer.crystallineColength_le_sum_dualDefect, New.Selmer.hasJumps_localLattice_summandRep, New.Selmer.sum_index_mul_defectBound_le, New.Selmer.eventually_exists_inClass_sum_dualDefect_le, New.Selmer.pow_chosenGroupOrder_eq_one_of_mem_inertiaAbove, New.Selmer.exists_tame_relations}
\lean{New.Selmer.eventually_exists_map_eq, New.Selmer.eventually_exists_map_eq_and_eval_one, New.Selmer.IsTameGeneratorFamily.dvd_eval_cyclotomicInv_sub_one, New.Selmer.normLength_addVal_pow, New.Selmer.ellVal_add_le_of_addVal_le, New.Selmer.eventually_exists_twistedLattice_and_block_lattices, New.Selmer.IsLatticeIn.det_one_sub, New.Selmer.exists_length_le_loss_add_crystallineSelmerLength, New.Selmer.ribet_columns}

\begin{lemma}[Lattices, tame relations and the accounting]\label{lem:ribet-inputs}
Let $V$, $c$ be a lift setting for $\vartheta$, $\psi$ an auxiliary twist of modulus $Q$ and ramification at most $\mathrm{ram}$, and $N=n_{\rm ch}$.
\begin{enumerate}[label=\textup{(\alph*)}]
\item If $u'$ has nondegenerate trace form and the trace of $\psi(R)$ plus the characters $O(-e)$, $e<t$, and the twisted lattice $\psi'R_0$ spans all matrices, then $u'$ has a lattice that is block diagonal with $\psi'R_0$ and the characters.
\item There are places $q_k\nmid\ell$, covering $Q$ and the places above $p_0p_1$ prime to $\ell$ up to conjugation, with tame generators $\delta_k$, and polynomials $o_k$ with $o_k(\delta)(\delta-1)=0$ on every $u_x$ for $\delta$ among the $\delta_k$ and the inverses of their
conjugates, such that the valuations of the $o_k(1)$ add up to at most $2\,\mathrm{ram}+8f\,v_\ell(N)$.
\item With divisible coefficients, the cocycles of the main summand $T$ at the lift that obey the relations and vanish at the $\tau_j$, the inverses of their conjugates and $\gamma$ have length at most $\mathrm{loss}_T+\mathrm{sel}_T$, where the losses of $T$ and the
crystalline colengths of the characters add up, normalised, to at most $r'\,\mathrm{loc}_\ell$.
\item Let $h\le v_\ell(N)$, $\mathrm{sep}=h_\Delta h$, $(k+n_{\rm test})\mathrm{sep}\le h_\Delta^{22}h$, $n_{\rm det}\le2(k+n_{\rm test})$, $\mathrm{tame}\le2\,\mathrm{ram}+8fv_\ell(N)$, $t'\le\mathrm t$, and let $v(d)$ be at most the losses, $\mathrm{sel}_T$, the odd class lengths $\mathrm{odd}_e\le v(X)$, $t'\,\mathrm{tame}$ and
$(2t'n_{\rm det}+k+1)\mathrm{sep}$. If the normalised losses add up to at most $\mathrm{loc}$ and $\mathrm{sel}_T\le\mathrm{sel}$, then $v(d)+2v_\ell(N)\le\mathrm{sel}+2t'\mathrm{ram}+\mathrm{loc}+e_\ell$.
\end{enumerate}
\end{lemma}

\begin{proof}
(a) Let $\rho$ be the block diagonal sum of $\psi'R_0$ and the characters, over $\bar\Q_\ell$. The traces of $u'$ and $\rho$ agree, the trace of a block sum being the sum of the traces; as both have nondegenerate trace forms on the spans of
their images, $u'$ is conjugate to $\rho$, which has entries in $O$.

(b) For $k\in Q$ take $o_k=\psi(\delta_k)-X$: on inertia above $Q$, $(g-\psi(g))(g-1)=0$, and $o_k(1)=\psi(\delta_k)-1$, of total valuation at most $\mathrm{ram}$ for each of the two copies. For a place above $p_0p_1$ prime to $\ell$ take
$o_k=(X^N-1)/(X-1)$: inertia there acts on a representation of the type through roots of unity of order dividing some $p_i^d-1$, which divides $N$ (\cref{lem:lift-settings}(b)), so $X^N-1$ kills it; $o_k(1)=N$. There are at most $4f$ such
places (\cref{lem:avatars}(a)) and two copies of each, which gives $8f\,v_\ell(N)$.

(c) A cocycle that kills the tame generators is unramified at the new places (\cref{lem:even-parts}(b)), so the members obeying the crystalline relations above $\ell$ are Selmer cocycles, of length at most $\mathrm{sel}_T$ modulo $\ell^M$ for every
$M$; the others are counted, modulo those, by the crystalline colength of $T$, which is bounded by its loss, the sum of the defects above $\ell$ (\cref{thm:local-loss}). Both bounds hold modulo every $\ell^M$, so by \cref{lem:divisible-cocycles} they hold
with divisible coefficients.

(d) $t'\,\mathrm{tame}\le2t'\mathrm{ram}+8f\mathrm t\,v_\ell(N)$; $\sum_e\mathrm{odd}_e\le\mathrm t\,v(X)$; and the multiples of $\mathrm{sep}$ are at most $4t'(k+n_{\rm test})\mathrm{sep}+(k+n_{\rm test})\mathrm{sep}+\mathrm{sep}\le((4\mathrm t+1)h_\Delta^{22}+h_\Delta)v_\ell(N)$, as $h\le v_\ell(N)$. Adding $2v_\ell(N)$,
the definition of $e_\ell$ gives the claim; $v$ is normalised so that $v(\ell)=1$.
\end{proof}

\begin{theorem}[Ribet's element for conjugate self-dual representations]\label{thm:ribet-columns}
Let $V$, $c$ be a lift setting for $\vartheta$ with $r'=r_V$, $t'=t_V$, and $\psi$ an auxiliary twist with ramification at most $\mathrm{ram}$. Let the $w_j$ form a triangle of primes with exponents $0<\mathrm{ty}(j)<t'$, up to
$\ell^{\mathrm r^8h}$, $h\le v(N)$. Let $u_x$, $x\in Y$, be finitely many conjugate self-dual representations of type $c$ and level $Q$, the $w_j$ and their conjugates, with $(g-\psi(g))(g-1)=0$ on inertia above $Q$; for $x\in\mathcal X$ let
$u_x$ have the trace of $\psi(R_x)$ plus $\mathrm{cyc}^{-e}$, $e<t'$, with $R_x$ a main summand and $\psi(R_x)$ having a lattice on which some $\gamma_x$ acts by a scalar. Let $S$, $\mathcal J$ be a trace ideal for $\ell^\nu$ and the relations
$\tau_j-1-X_j(\varphi_j-1)H_j(\varphi_j)$, $H_j(1)=1$, and let the Selmer group of $\psi(R_\pi)$ with divisible coefficients have length at most $\mathrm{sel}<\infty$. Then the $\det(1-u_x(g))$ lie in $S$ and vanish on $\mathcal X$, and some $f\in S$
multiplies each of them into $\mathcal J$, has on $\mathcal X$ no monomial properly dividing $\prod X_j$, and its coefficient $d$ of $\prod X_j$ at $\pi$ satisfies
\[
  v(d)+2v(N)\le\mathrm{sel}+2t'\,\mathrm{ram}+r'\,\mathrm{loc}_\ell+e_\ell.
\]
\end{theorem}

\begin{proof}
\emph{The coefficient ring.} Choose a coefficient ring $O$ that carries all the data: the finitely many polynomials $H_j$ over the integers of $\bar\Q_\ell$, with $H_j(1)=1$, are defined over every large $O$, and so are lattices in the $u_x$ and models
$R_0$, $\psi'$, $\zeta'$ of the main summands with $\zeta'-1$ a unit; at $\pi$ take the model with the Selmer bound. By \cref{lem:ribet-inputs}(a) the lattices may be taken block diagonal on $\mathcal X$. \emph{Doubling.} With the inverses of the
conjugates of the $\tau_j$, $\varphi_j$ added, both copies using $X_j$ and $H_j$, $\mathcal J$ is again a trace ideal: on the copy of $(\tau_j^c)^{-1}$, conjugate self-duality turns each trace into a power of $\mathrm{cyc}^{-1}$ times a trace at $g'=(g^c)^{-1}$, and as
$\mathrm{cyc}^{-1}$ is $1$ at $\tau_j$ and $\equiv1$ modulo $\ell^\nu$ at $\varphi_j$, the new relation is a unit times the old one at $g'$, up to multiples of $\ell^\nu$ times traces. The factors $\varepsilon_j=H_j$ have the value $\equiv1$ modulo $\ell^\nu$ at the scalars
$\mathrm{cyc}^{-e}(\varphi_j)$ by which the $\varphi_j$ act on the characters.
\emph{The hypotheses of \cref{thm:ribet-lemma}.} (i) and (ii): on $\mathcal X$ the lattice is the block $\psi'R_0$, which spans all matrices (\cref{prop:main-summands}), plus the characters $\mathrm{cyc}^{-e}$ with $\xi_0=1$; $\gamma_x$ acts on the block by $\zeta'$ and is an
inertia element off $\ell$, where the cyclotomic character is trivial. (iii): the element $g_0$ of the triangle separates the characters, $\prod_{k\ne e}(u^e-u^k)\mid\ell^{h_\Delta h}=C$, and the types are monotone. (iv): \cref{lem:ribet-inputs}(c). (v): the elements
$\delta$ of \cref{lem:ribet-inputs}(b), trivial on the characters (the cyclotomic character is trivial on inertia off $\ell$). (vi): \cref{thm:triangle-cost} with the odd classes bounded by the odd class bound and the crystalline colengths of the
characters by \cref{lem:ribet-inputs}(c); the detection on the test points of the triangle and $g_0$, together with the tame generators, which a character cocycle vanishing there and unramified off $\ell$ must respect.
\emph{The determinants.} $\det(1-u_x(g))$ is the determinant of the combination $\sum_{g'}\omega_{g'}(u_x(g')-1)$ with $\omega=-1$ at $g$ and $0$ elsewhere, computed on the lattice. \emph{The bound.} By \cref{thm:ribet-lemma} the coefficient comes from $c\in O$
with $v(c)$ at most $b_0+\sum_eb_e+t'(\sum_kv(o_k(1))+2\#s\,v(C))+(\#J+1)v(C)$, and \cref{lem:ribet-inputs}(d), with $k=\#J$, $n_{\rm test}$ the number of test points and $g_0$, $n_{\rm det}=\#s$, gives the stated inequality.
\end{proof}

\subsection{The bound for denominators}\label{ss:denominator-bound}
A \emph{lift generator} of a Hecke-stable space $Z$ of forms on $V$ is a form of $Z$ in the range of the projection of a type in lift setting for $\vartheta$, lifting a Hecke eigensystem that has a main summand whose twist by
$\vartheta^{-1}_\ell$ has a Selmer group with divisible coefficients of length at most $\mathrm{sel}$.
\lean{New.Selmer.exists_isUnit_cyclotomicChar_pow_sub_one, New.Selmer.exists_cyclotomicChar_pow_sub_one_dvd_pow_padicValNat, New.Selmer.exists_cyclotomicChar_pow_sub_one_dvd, New.Selmer.IsWeakLift, New.Selmer.IsNonzeroIn, New.Selmer.IsLiftWithTwistedLattice, New.Selmer.liftValuationBound, New.Selmer.invAvatar_eq_one_of_mem_inertiaAbove, New.Selmer.exists_eventually_hasRootChain_mostFrobeniiAt, New.Selmer.exists_taylorWilesPlaces_hasRootChain, New.Selmer.norm_charpoly_coeff_le_one_of_continuous, New.Selmer.exists_det_one_sub_ne_zero, New.Selmer.exists_prod_sub_det_one_sub_eq_zero, New.Selmer.exists_dvd_prime_pow_succ, New.Selmer.normalizingAvatar_tameGenerator, New.Selmer.exists_isCommonRibetElement_ellVal_le, New.Selmer.exists_ellVal_le_and_diagOp_mem_ellIntegralForms, New.Selmer.exists_heckePolyOnForms_eq_zero, New.Selmer.typeProjection_sub_starProjection_mem, New.Selmer.IsLiftGenerator, New.Selmer.IsLiftGenerator.mono, New.Selmer.span_mem_invtSubmodule_heckePolyOnForms, New.Selmer.IsLiftGenerator.typeProjection_eq_zero_or_mem_range, New.Selmer.IsLiftGenerator.exists_isMainSummandLike, New.Selmer.span_isLiftGenerator_mem_invtSubmodule_heckePolyOnForms, New.Selmer.span_isLiftGenerator_mem_invtSubmodule_typeProjection, New.Selmer.span_isLiftGenerator_mem_invtSubmodule, New.Selmer.typeProjection_sub_starProjection_mem_of_isOrthogonalToFalseLifts, New.Selmer.mem_integralForms_ofIso_iff, New.Selmer.IsLiftGenerator.exists_isMainSummandLike_hasTwistedLattice, New.Selmer.diagOp_toSmoothForm_typeProjection, New.Selmer.diagOp_typeProjection_eq_zero_of_mem_span, New.Selmer.exists_linearMap_eq_smul_mem_integralForms, New.Selmer.PreservesHeckeEigenspaces, New.Selmer.preservesHeckeEigenspaces_toLin_diagonal, New.Selmer.exists_operator_ellVal_le_liftValuationBound, New.Selmer.exp_le_norm_det_smul_of_ellVal_le, New.Selmer.exists_operator_det_clearingBound, New.Selmer.denominatorLe_clearingBound}

\begin{lemma}[Separation, weak lifts and lift generators]\label{lem:denominator-inputs}
\begin{enumerate}[label=\textup{(\alph*)}]
\item There are $h\le v_\ell(N)$, with $h=0$ unless $\ell\le h_\Delta^2$, and a topological generator $g_0$ of $\Gal(\Omega/E)$ modulo the kernel of the cyclotomic character with $\mathrm{cyc}(g_0)^a-1\mid\ell^h$ for $0<a<\mathrm t$.
\item Some split place $w$ off the places of $O$ and their conjugates has the property that $H_w(X_*)$ kills every nonzero form of base level in a weak lift.
\item If $e_c$ is the type projection, $Z$ is stable under it and under the Hecke polynomials, $y$ is orthogonal to the false lifts, and every nonzero isotypic form in the range of $e_c$ outside $Z$ belongs to a weak lift, then
$e_c(y-\mathrm{pr}_Zy)$ lies in the span of the weak lifts of base level.
\item The span of the lift generators is stable under the Hecke polynomials and the type projections; a lift generator either is killed by $e_c$ or is a nonzero form of a system in its range, and every system containing a nonzero lift
generator lifts a main summand with a twisted lattice.
\end{enumerate}
\end{lemma}

\begin{proof}
(a) If $\ell>h_\Delta^2$, take $\gamma$ in an inertia group of $\Q$ at $\ell$ whose cyclotomic character reduces to a generator $b$ of $(\Z/\ell)^\times$; it fixes $E$, which lies in a cyclotomic field of conductor prime to $\ell$. As $b$ has order $\ell-1$ and
$0<a<\mathrm t\le h_\Delta^2<\ell$, $b^a\ne1$, so $u^a-1$ is a unit and $h=0$. If $\ell\le h_\Delta^2$, take a Frobenius element at a place of norm $p_i\ne\ell$; then $u^a-1=p_i^a-1$ divides $N$ (\cref{lem:lift-settings}(b)), and a nonzero $\ell$-adic
integer is a unit times $\ell^{v_\ell}$, so $u^a-1\mid\ell^{v_\ell(N)}$. In both cases a topological generator modulo the kernel may replace the element, the divisibility passing to it.

(b) For a split $w$ off the level, a nonzero form of base level in a weak lift has the roots $X_0q^i$, $i<t$, in its Hecke polynomial (\cref{lem:root-chains}(b)). The last of them is $X_*^{-1}$, so the reversed polynomial, by which $H_w$ acts, vanishes at $X_*$.

(c) Let $U'$ be the part of the range of $e_c$ orthogonal to $Z$, and $K'$ the forms of $U'$ killed by $H_w(X_*)$ for the place $w$ of (b); they are false lifts, being in $Z^\perp$, in the image of $e_c$, and killed by the Hecke polynomial at the lift
root. An isotypic $v\in U'$ outside the span $M$ of the weak lifts lies in $K'$ by the hypothesis and (b), hence is orthogonal to itself; so $U'\cap K'^\perp\subseteq M$ (\cref{lem:root-chains}(d), the Hecke operators at $w$ and at other places
commuting). The form $e_c(y-\mathrm{pr}_Zy)$ lies in $U'$, and in $K'^\perp$: $e_c$ is self-adjoint (\cref{lem:lift-settings}(c)) and fixes $K'$, and both $y$ and $\mathrm{pr}_Zy$ are orthogonal to the false lifts at $w$, the latter because $\mathrm{pr}_Zy\in Z\perp K'$.

(d) A lift is an eigenvector of every Hecke polynomial, so the span of lifts is stable. The type projections are central idempotents (\cref{lem:lift-settings}(c)) and a lift generator lies in the range of one of them inside the image of
$e_0$; idempotents of distinct tuples are orthogonal, so $e_c$ either kills it or fixes it, and a nonzero lift generator is a nonzero form of the system of its Hecke eigenvalues. Such a system lifts a main summand $R$ with the trace
identity of \cref{thm:main-summands-exist}, and a twisted lattice exists by the scalar of large order of \cref{def:main-summand}; for the generator $v_0$ with the Selmer bound the lattice can be taken with that bound.
\end{proof}

\begin{theorem}[Bound for denominators]\label{thm:denominator-bound}
\begin{enumerate}[label=\textup{(\alph*)}]
\item Let $\vartheta$ be auxiliary with ramification at most $\mathrm{ram}$ and let $Z$ be spanned by lift generators. Then an operator $\Lambda$ on $Z$ preserves the Hecke eigenspaces at the split places off the level, has eigenvalues of
valuation at most $\mathrm{sel}+2t_O\,\mathrm{ram}+r_O\,\mathrm{loc}_\ell+e_\ell$, and takes the projections to $Z$ of integral forms orthogonal to the false lifts to integral forms.
\item If $\mathrm{ram}$ and $\mathrm{sel}$ are at most $e'_\ell$, some operator on $Z$ has determinant of norm at least $\exp(-\dim Z\cdot r_O\beta_\ell\log\ell)$ and takes the projection of every integral form orthogonal to the false lifts to an
integral form whose pairings with integral forms are integral; that is, the projection to $Z$ has denominator at most $\ell^{\beta_\ell}$.
\end{enumerate}
\end{theorem}

\begin{proof}
(a) \emph{One generator.} Let $v_0\ne0$ be a lift generator with the Selmer bound, of type $c$ and system $x_0$. Base data $B$ for $O$, $c$ and $\vartheta$ exist (\cref{thm:attached}(c)): the forms of type $c$ are forms of $B$ of base level, and the
normalised representations of its systems at every auxiliary level are of type $c$, with the quadratic relation above $Q$. \emph{The weak-lift selector.} The finitely many base systems $x$ that are not weak lifts have elements $g_k$ with
$\det(1-u_x(g_k))=\mathrm{value}_k\ne0$, integral as $u_x$ is continuous on a compact group; put $\delta_0=\prod_k(\mathrm{value}_k-\delta(g_k))$ with $\delta(g)=(x\mapsto\det(1-u_x(g)))$ and $n_0=\prod_k\mathrm{value}_k$. Then $\delta_0$ vanishes at the systems that are not weak lifts,
and a base system with $\delta_0\ne0$ is a weak lift. \emph{The auxiliary places.} With $g_0$, $h$ of \cref{lem:denominator-inputs}(a), \cref{thm:triangles} gives a triangle of Taylor--Wiles places of every level $\ell^{\nu}$, with Frobenius data in the
set where, by \cref{lem:root-chains}(c), every weak lift of base level congruent to $x_0$ has a root chain with an integral divisor $e_j$, $e_j(1)=1$; the set contains most Frobenii near $\delta^{\ell^\nu}g$ for every $g$, as the twist of the main
summand by a power of the scalar inertial element has $\norm{\det(1-\cdot)}\ge1$ near suitable points. Take $\nu$ with $2n!(n_0n_{\rm Aut}^2)^{\#J+1}\mid\ell^\nu$. \emph{Ribet's element.} By \cref{thm:ribet-columns} applied to the systems of the auxiliary level,
every trace ideal has a common Ribet element for the $\delta(g)$ with $v(N^2d)$ at most $L=\mathrm{sel}+2t_O\mathrm{ram}+r_O\mathrm{loc}_\ell+e_\ell$. Reading polynomials over $\mathbb T$ as families of polynomials, one for each system, the Ribet element becomes $\mathcal D\in\mathbb T[\alpha]$
multiplying every $\delta(g)$ into the tied ideal $\mathfrak a$; as $\delta_0-n_0$ lies in the ideal of the $\delta(g)$, $(\delta_0-n_0)\mathcal D\in\mathfrak a$, its coefficients below $\prod\alpha_j$ vanish on $\mathcal X$, and its top coefficient at $x_0$ has the bound. \emph{The
clearing.} Let $l$ be a form of base level with $N l$ integral, $y$ supported on $\mathcal X$ and agreeing with $l$ where $\delta_0\ne0$, and $A$ the indicator of the systems congruent to $x_0$ (\cref{prop:dictionary}(a)), with $A^{\#J+1}=A$. By
\cref{thm:deformed-norm}(d) with $d=N$, $N^2A\mathcal D_Jy$ is integral, the common multiple $n_{\rm Aut}$ of the stabiliser orders of these base data being $N$; and $s=N^2\mathcal D_J$ has $v(s(x_0))\le L$. \emph{The operator.} Let $G$ be $s$ composed with $e_c$, scaling each system. For $y$ integral and orthogonal to the false lifts,
$e_c(\mathrm{pr}_Zy)$ is supported on the systems of lift generators, and agrees with $e_c(y)$ at the weak lifts by \cref{lem:denominator-inputs}(c) and (b); so $G(\mathrm{pr}_Zy)$ is integral. Every lift generator is an eigenvector of $G$, with eigenvalue of valuation at
most $L$ at $v_0$. \emph{All generators.} Choose a basis of $Z$ among the lift generators; the operators for the basis vectors are diagonal in it, and by \cref{lem:root-chains}(a) one integral combination $\Lambda$ has at every basis vector an eigenvalue at least as
large in norm, of valuation at most $L$. A diagonal operator in a basis of Hecke eigenvectors preserves the Hecke eigenspaces, and integral combinations preserve the integrality property.

(b) If $Z=0$ the identity serves. Otherwise the lift setting of a generator bounds $r_O\le\mathrm r$ and $t_O\le\mathrm t$, and $N\Lambda$, with $\Lambda$ from (a), has eigenvalues of valuation at most $L+v_\ell(N)\le r_O\beta_\ell$, since
$\mathrm{sel}+2t_O\mathrm{ram}+r_O\mathrm{loc}_\ell+e_\ell+v_\ell(N)\le r_O\beta_\ell$ for $\mathrm{sel},\mathrm{ram}\le e'_\ell$ and $1\le r_O$ ($\beta_\ell=\mathrm{loc}_\ell+(1+2\mathrm t)e'_\ell+e_\ell+v_\ell(N)$). By \cref{lem:root-chains}(a) its determinant has norm at least $\ell^{-\dim Z\,r_O\beta_\ell}$. The image of an integral form under
$N\Lambda$ is integral, and its pairings with forms integral at all places above $\ell$ are integral, the Gram matrix being integral and $N$ clearing the stabiliser orders.
\end{proof}
